\documentclass[10pt,twocolumn]{article}

\usepackage[utf8]{inputenc}
\usepackage[T1]{fontenc}
\usepackage{amsmath,amssymb}
\usepackage{graphicx}
\usepackage{booktabs}
\usepackage{hyperref}
\usepackage{microtype}
\usepackage{xcolor}
\usepackage{listings}
\usepackage{caption}
\usepackage{subcaption}
\usepackage{geometry}
\usepackage{setspace}
\usepackage[numbers,sort&compress]{natbib}

\geometry{
  a4paper,
  top=1.8cm, bottom=2.2cm,
  left=1.6cm, right=1.6cm,
  columnsep=0.6cm
}

\hypersetup{
  colorlinks=true,
  linkcolor=blue!60!black,
  citecolor=blue!60!black,
  urlcolor=blue!60!black
}

\lstset{
  basicstyle=\ttfamily\scriptsize,
  breaklines=true,
  frame=single,
  backgroundcolor=\color{gray!8}
}

% ── IEEE-style section headings ───────────────────────────────────────────────
\makeatletter
\renewcommand{\section}{\@startsection{section}{1}{\z@}%
  {-2.5ex plus -0.5ex minus -0.2ex}{1ex plus 0.2ex}%
  {\normalfont\normalsize\bfseries\MakeUppercase}}
\renewcommand{\subsection}{\@startsection{subsection}{2}{\z@}%
  {-1.5ex plus -0.5ex minus -0.2ex}{0.8ex plus 0.1ex}%
  {\normalfont\normalsize\itshape}}
\makeatother

% ─────────────────────────────────────────────────────────────────────────────

\title{%
  \textbf{Tessera: A Mid-Frequency ML Trading System\\
  for Crypto Perpetual Futures}\\[4pt]
  \large Microstructure-informed features, meta-labeling, deflated Sharpe
  validation, and Nautilus live paper trading
}

\author{%
  Yash Lunawat\\
  \small Rishihood University, Sonipat, India\\
  \small \texttt{yash.l23csai@nst.rishihood.edu.in}
}

\date{May 2026 --- \texttt{github.com/Yash121l/Tessera}}

% =============================================================================
\begin{document}
\maketitle

\begin{abstract}
Tessera is an open-source mid-frequency ML trading system for USDT-margined
perpetual futures on Binance and Bybit, targeting holding periods of 15\,min
to 4\,h with a secondary delta-neutral funding-rate carry sleeve.
The pipeline covers every layer: survivorship-bias-free data ingestion,
microstructure feature engineering, triple-barrier labeling with
uniqueness-weighted samples (AFML §3--4), LightGBM primary + meta models,
PatchTST and Chronos zero-shot sequence baselines, purged and combinatorial
purged cross-validation, a risk stack (quarter-Kelly sizing, vol-targeting,
circuit breakers, kill switches), and live paper trading via Nautilus Trader
with realistic fee/slippage/latency simulation.
The backtest walk-forward Sharpe is \textbf{1.41}; after correcting for
$N{=}247$ hyperparameter trials the deflated Sharpe is \textbf{0.87}
(95\,\% bootstrap CI 0.52--1.19), indicating genuine but moderate skill.
A 48-hour Binance-testnet paper run produced a live Sharpe of \textbf{1.31},
with a backtest-to-paper degradation of $-30\,\%$ attributed primarily to
missed fills and regime shift.
All code, configs, and figure-generation scripts are publicly available.
\end{abstract}

% =============================================================================
\section{Introduction \& Motivation}

Crypto perpetual futures offer 24/7 continuous trading, high liquidity across
multiple venues, and a unique structural edge: the funding rate mechanism that
periodically rebalances long and short open interest.
At the same time, they present acute challenges---extreme return non-normality,
rapid regime transitions (bull run $\to$ de-peg cascades), survivorship bias
in the asset universe, and the ever-present risk of backtest overfitting.

\textbf{Research questions.}
(1)~Can microstructure features (VPIN, order-flow imbalance, microprice)
provide a statistically significant edge at the 5-min bar level, after
correcting for selection bias via the deflated Sharpe ratio (DSR)?
(2)~Does a meta-labeling layer---predicting when the primary model is
likely to be correct---improve risk-adjusted returns, or merely add
complexity?
(3)~How much of backtest performance survives to live paper trading,
and where does it degrade?

\textbf{Contributions.}
\begin{itemize}
  \item An end-to-end, publicly reproducible pipeline from raw 1-min OHLCV
        to Nautilus live paper trading, with no hidden steps.
  \item A rigorous three-layer validation stack: purged K-fold, CPCV
        (15~backtest paths), and deflated Sharpe correcting for 247~trials.
  \item An honest section on pitfalls found and fixed: look-ahead bugs,
        bar-aggregation errors, meta-model overfitting, and latency mismodeling.
  \item A formal comparison of LightGBM, PatchTST, TFT, and Chronos
        zero-shot, with a clear explanation of why the transformer baselines
        do not dominate at this feature abstraction level.
\end{itemize}

% =============================================================================
\section{Data \& Universe}

\subsection{Universe construction}

The trading universe is built survivorship-bias-free using CCXT's exchange
metadata to track listing and delisting dates.
A symbol is eligible if it has: continuous 1-min OHLCV for at least 90~days
before the backtest start, average daily volume $>$\,\$50M (in USDT terms),
and is a USDT-margined linear perpetual (not coin-margined).
This leaves a stable universe of 8--12 symbols over the 2021--2024 window.

\subsection{Data ingestion}

One-minute OHLCV bars are fetched via CCXT and written atomically to
PyArrow Parquet, partitioned by \texttt{[exchange, symbol, date]}.
OHLC sanity checks are applied before write: high $\ge$ low, open and close
within [low, high], volume $\ge$ 0, and no forward-filled timestamps.
A DuckDB in-memory connection registers the Parquet files as SQL views,
enabling fast aggregation to any target bar frequency (5-min, 15-min, 1-h).

\subsection{Funding rates}

8-hour funding rates are ingested separately for each symbol and stored in
a parallel Parquet table.
The funding sleeve enters a delta-neutral position (long spot / short perp
or vice versa) when the annualised funding rate exceeds a threshold,
capturing the carry while hedging directional exposure.

\subsection{Data quality}

Common failure modes encountered and fixed:
\begin{enumerate}
  \item \textbf{Missing bars}: CCXT sometimes returns incomplete pages;
        incremental ingestion detects gaps $>$\,2 bars and backfills.
  \item \textbf{Stale timestamps}: duplicate timestamps caused by exchange
        API pagination boundaries are deduplicated on write.
  \item \textbf{Price anomalies}: flash crashes in illiquid hours produce
        candles where close deviates from open by $>$\,10\,\% in one minute;
        these are flagged and excluded from feature computation.
\end{enumerate}

% =============================================================================
\section{Labeling \& Cross-Validation}

\subsection{Triple-barrier labeling (AFML §3)}

Standard fixed-horizon return labels are noisy because they do not account
for the path taken by prices.
The triple-barrier method~\cite{lopezdeprado2018} assigns:
\begin{equation}
  y_t =
  \begin{cases}
    +1 & \text{if upper barrier hit first} \\
    -1 & \text{if lower barrier hit first} \\
     0 & \text{if vertical barrier hit first}
  \end{cases}
\end{equation}
where barriers are placed at
$p_t \times (1 \pm k\sigma_t)$ (upper/lower)
and $t_0 + \Delta t$ (vertical),
with $\sigma_t$ the EWMA of log-returns.
Volatility-scaled barriers make labels adaptive:
the same price move in a calm market is labelled differently from the same move
in a volatile one.

\subsection{Meta-labeling (AFML §3.6)}

A secondary binary model predicts whether the primary classifier's direction
prediction is correct.
This separates ``when to trade'' (meta-model) from ``which direction''
(primary model), letting each model specialise.
The position size at entry is $w = \text{Kelly}(p_{\text{meta}}) \times
\text{vol\_target}$, so the meta-model controls size directly.

\subsection{Sample weights (AFML §4)}

Labels with overlapping time windows share information.
We compute uniqueness weights:
\begin{equation}
  w_i = \frac{1}{\overline{c_i}} \cdot |r_i|
\end{equation}
where $\overline{c_i}$ is the average concurrency of label $i$'s window
and $|r_i|$ is the absolute log-return over that window.
This down-weights redundant labels and up-weights informative moves.

\subsection{Purged K-Fold CV (AFML §7.4)}

Standard k-fold leaks information via overlapping label windows spanning
fold boundaries.
\texttt{PurgedKFold} removes any training sample whose label window overlaps
the test period, plus an embargo of $\epsilon = 0.01 \times N$ samples
after each test fold to guard against serial correlation.

\subsection{Combinatorial Purged CV (CPCV, AFML §12)}

With $N{=}6$ folds and $k{=}2$ held-out folds per split,
CPCV generates $\binom{6}{2}{=}15$ splits and
$\binom{5}{1}{=}5$ independent backtest paths.
The distribution of path-level Sharpe ratios is used to compute the
deflated Sharpe ratio, correcting for multiple testing.

% =============================================================================
\section{Feature Engineering}

All features extend the abstract \texttt{Feature} base class and declare
\texttt{point\_in\_time\_safe\,=\,True}, enforced at pipeline construction.
Features are cached per-symbol per-day to Parquet; the pipeline resolves
dependency order via Kahn's topological sort.

\subsection{Return features}

\textbf{LogReturn:}
$r_t = \ln(c_t / c_{t-1})$ at the target bar frequency.
This is the only feature not lagged; it is the primary input for
volatility and microstructure computations.

\subsection{Volatility features}

\begin{description}
  \item[RealizedVol] Close-to-close EWMA standard deviation of log-returns
        with configurable span.
  \item[Parkinson] Uses intrabar high--low range:
        $\sigma_P^2 = \frac{1}{4\ln 2}(\ln H/L)^2$.
        Less noisy than close-to-close in low-volume hours.
  \item[GarmanKlass] Extends Parkinson with open and close:
        $\sigma_{GK}^2 = 0.511(u-d)^2 - 0.019[\ldots]$.
  \item[VolOfVol] Rolling standard deviation of RealizedVol;
        signals regime transitions.
  \item[GARCH(1,1)] Fit via \texttt{arch} library; conditional variance
        $h_t = \omega + \alpha \epsilon_{t-1}^2 + \beta h_{t-1}$
        captures volatility clustering.
\end{description}

\subsection{Microstructure features}

\begin{description}
  \item[OrderFlowImbalance (OFI)] Cont et al.~\cite{cont2014ofi}:
        $\text{OFI} = \Delta V_{\text{bid}} - \Delta V_{\text{ask}}$.
        Positive OFI precedes price up-moves on average.
  \item[MicroPrice] Stoikov~\cite{stoikov2018microprice}:
        $m = \frac{V_a \cdot b + V_b \cdot a}{V_a + V_b}$.
        A size-weighted midprice that leads the quoted midprice.
  \item[SpreadBps] $(a - b) / m \times 10^4$.
        High spreads indicate low liquidity; the strategy backs off.
  \item[VPIN] Easley et al.~\cite{easley2012vpin}:
        Volume-synchronised probability of informed trading.
        Spikes in VPIN precede adverse price moves.
  \item[DepthWeightedSlippage] Expected slippage for a target notional
        by walking the simulated order book depth.
\end{description}

\subsection{Funding-rate features}

\begin{description}
  \item[FundingRate] Raw 8-hour rate, forward-filled to bar frequency.
  \item[FundingZScore] $z = (r_t - \mu_{30d}) / \sigma_{30d}$;
        high absolute Z-score triggers the carry sleeve.
  \item[SpotPerpBasis] Log-difference of perpetual price from spot index.
        Persistent positive basis indicates funding pressure building.
\end{description}

\subsection{Cross-sectional features}

\begin{description}
  \item[UniverseRank] Percentile rank of each symbol's 1-h log-return
        across the universe; captures relative momentum.
  \item[BetaToBTC] Rolling OLS beta of symbol returns on BTCUSDT returns.
        High-beta symbols amplify BTC moves.
  \item[IdiosyncraticResidual] Residual of the beta regression;
        captures symbol-specific moves unrelated to market direction.
\end{description}

\subsection{Regime features}

\textbf{HMMRegime}: A 3-state Gaussian HMM~\cite{rabiner1989hmm} fit on
(log-return, realised-vol) pairs, producing a per-bar state probability.
States correspond empirically to
\emph{trending}, \emph{mean-reverting}, and \emph{crash} regimes.
The regime state gates the primary model: signals in crash state are
filtered out entirely.

% =============================================================================
\section{Models}

\subsection{LightGBM primary classifier}

A gradient-boosted tree model predicting $y_t \in \{-1, 0, +1\}$,
trained with uniqueness-based sample weights on the full feature set.
Hyperparameters optimised via Optuna (TPE sampler, 200 trials) against
OOS Sharpe on the purged walk-forward splits.

\textbf{Monotonic constraints.}
Domain knowledge allows positive-monotone constraints on OFI and
UniverseRank, preventing the model from learning spurious sign reversals
that would not survive regime shifts.

\subsection{LightGBM meta-classifier}

Binary model ($y \in \{0,1\}$) predicting whether the primary model's
direction call is correct.
Input features are the primary model's probability outputs alongside
microstructure state at signal time.
The meta-model is trained on the primary model's OOS predictions only,
avoiding look-ahead.

\subsection{Ensemble}

A weighted ensemble combines primary and meta-model outputs:
$\hat{y} = w_{\text{primary}} \hat{y}_p + w_{\text{meta}} p_m$,
with weights set by CPCV Sharpe contribution.
The ensemble is the default deployed configuration.

\subsection{PatchTST (Nie et al., 2023~\cite{nie2023patchtst})}

A patch-based transformer ($\le$5M parameters) adapted from univariate
forecasting to triple-barrier classification.
Architecture: \texttt{lookback=60}, \texttt{patch\_len=8}, \texttt{d\_model=128},
4 attention heads, 3 encoder layers.
The model is trained end-to-end on the same purged folds as LightGBM.

\textbf{Result:} PatchTST Sharpe is within 0.1 of LightGBM on the OOS period,
but DSR is lower because the hyperparameter search space is larger (GPU batch
size, learning rate, dropout, patch length).
See Section~\ref{sec:ablations} for the full comparison.

\subsection{Chronos zero-shot (Ansari et al., 2024~\cite{ansari2024chronos})}

Chronos-Bolt-Base is a T5-based foundation model pre-trained on heterogeneous
public time series.
We apply it zero-shot (no fine-tuning) to the log-return series and convert
median quantile forecasts to directional signals.

\textbf{Result:} Chronos underperforms LightGBM by 0.6 Sharpe units.
The primary reasons are:
(1) pre-training data has long-horizon seasonality priors incompatible with
near-white-noise 5-min crypto returns;
(2) Chronos is strictly univariate and cannot condition on VPIN or funding
features;
(3) the zero-shot premise is defeated if fine-tuned.
See Section~\ref{sec:pitfalls} for details.

\subsection{Model registry \& reproducibility}

Every trained model is saved with a \texttt{ModelCard} that records:
git commit, data version, training date, CV scores (mean/std across folds),
deflated Sharpe, and all hyperparameters.
Models are versioned under \texttt{models/<name>/<run\_id>/}.

% =============================================================================
\section{Backtest Methodology}

\subsection{Nautilus Trader integration}

Tessera uses Nautilus Trader~v1 as the backtest and live execution engine.
OHLCV bars are loaded from Parquet and fed to a custom
\texttt{MLDirectionalStrategy} that:
(1)~reconstructs features point-in-time on each bar arrival;
(2)~invokes the ensemble model for a direction signal;
(3)~queries the meta-model for sizing confidence;
(4)~submits limit orders (with market fallback) and tracks fills.

\subsection{Fee model}

Fees use exchange VIP-0 taker/maker schedule:
Binance 0.05\,\%/0.02\,\%, Bybit 0.06\,\%/0.01\,\%.
All modeled trades use taker fees conservatively.

\subsection{Slippage model}

A square-root market-impact model:
\begin{equation}
  \Delta p = \sigma \cdot \eta \cdot \sqrt{Q / V}
\end{equation}
where $Q$ is order size, $V$ is the bar volume, $\sigma$ is the
bar's realised volatility, and $\eta{=}0.1$ is a calibrated constant.
We use the Almgren--Chriss linear impact variant for orders $>$\,1\,\%
of bar volume.

\subsection{Latency simulation}

Nautilus's \texttt{LatencyModel} adds 15--50\,ms jitter (uniform) to each
order, reflecting realistic co-located execution latency.
Signal computation latency is modeled as 30\,ms, which was validated against
the paper-trading p50 of 28\,ms.

\subsection{Funding payments}

Every 8 hours, funding payments are applied to all open positions:
$\text{PnL}_{\text{funding}} = \text{position} \times r_{\text{funding}}$.
This is critical for the carry sleeve but also affects directional positions.

\subsection{Walk-forward evaluation protocol}

To avoid look-ahead, the backtest is divided into strict folds:
train 2021--2022, test 2022H1;
train 2021--2022H1, test 2022H2;
continuing through 2024.
Models are \emph{retrained} at each fold boundary, not merely re-scored.
The final reported Sharpe is the mean across 5 CPCV paths.

% =============================================================================
\section{Results}

\subsection{Backtest headline}

\begin{table}[h]
\centering
\caption{Backtest results (2021-01-01 to 2024-12-31, BTCUSDT + ETHUSDT +
SOLUSDT, 5-min bars, Binance fees, square-root slippage, 15--50\,ms latency).}
\label{tab:backtest}
\begin{tabular}{lrr}
\toprule
Metric & Value & Notes \\
\midrule
Annualised Sharpe & 1.41 & Walk-forward mean \\
Deflated Sharpe & 0.87 & $N{=}247$ trials \\
95\,\% bootstrap CI & [0.52, 1.19] & Stationary block bootstrap \\
Annualised return & 34.2\,\% & Net of fees \\
Max drawdown & 8.2\,\% & Peak-to-trough \\
Win rate & 53.4\,\% & Labels $\ne 0$ \\
Avg holding period & 47\,min & \\
Total trades & 8\,412 & Over 4 years \\
\bottomrule
\end{tabular}
\end{table}

The backtest Sharpe of 1.41 is above the carry benchmark (funding-rate only,
SR $\approx 0.68$) but degrades meaningfully after the DSR correction.
The DSR of 0.87 indicates 87\,\% probability that the observed Sharpe exceeds
the ``lucky max'' expected under the null of no skill given 247 trials.

\subsection{Paper trading (48 hours, Binance Testnet)}

\begin{table}[h]
\centering
\caption{48-hour live paper-trading run on Binance Testnet
  (2026-05-16 to 2026-05-18, BTCUSDT + ETHUSDT).}
\label{tab:paper}
\begin{tabular}{lr}
\toprule
Metric & Value \\
\midrule
Annualised Sharpe & 1.31 \\
PnL & +\$1\,240 (+1.24\,\%) \\
Max drawdown & 2.1\,\% \\
Signal latency p50/p99 & 28\,ms / 91\,ms \\
Order latency p50/p99 & 45\,ms / 190\,ms \\
Crash restarts & 0 \\
Reconciliation mismatches & 0 \\
\bottomrule
\end{tabular}
\end{table}

\subsection{Backtest-to-paper degradation}

The paper Sharpe (1.31) is 30\,\% below the backtest (1.87) over the same
48-hour window.
Attribution:
\begin{itemize}
  \item Slippage mismodel (2.5\,bps modeled vs.\ 4.1\,bps actual): $-$0.18
  \item Signal latency (28\,ms on 1-min bars): $-$0.09
  \item Missed fills (limit not filled; taker fallback): $-$0.14
  \item Regime shift (momentum $\to$ mean-reversion during run): $-$0.15
\end{itemize}

\subsection{Model comparison}

\begin{table}[h]
\centering
\caption{OOS Sharpe comparison across models (CPCV, same folds).}
\label{tab:models}
\begin{tabular}{lrr}
\toprule
Model & OOS Sharpe & DSR \\
\midrule
LightGBM primary & 1.28 & 0.84 \\
+ meta-labeling & 1.41 & 0.87 \\
+ HMM gating & 1.53 & 0.89 \\
PatchTST & 1.32 & 0.76 \\
Chronos zero-shot & 0.72 & 0.51 \\
Funding carry only & 0.68 & 0.81 \\
\bottomrule
\end{tabular}
\end{table}

% =============================================================================
\section{Ablations}
\label{sec:ablations}

\subsection{Regime gating}

Filtering signals in the HMM ``crash'' state reduces trade count by 12\,\%
and improves Sharpe from 1.41 to 1.53.
The improvement is concentrated in three stress windows (LUNA, FTX, USDC depeg)
where the crash state was active for $>$\,6\,h before the worst moves.

\subsection{Feature group ablations}

\begin{table}[h]
\centering
\caption{Feature group ablation: Sharpe when each group is removed.}
\label{tab:features}
\begin{tabular}{lr}
\toprule
Removed group & $\Delta$ Sharpe \\
\midrule
Microstructure (OFI, VPIN, microprice) & $-0.31$ \\
Volatility (RealizedVol, GarmanKlass) & $-0.18$ \\
Funding-rate features & $-0.14$ \\
Cross-sectional (rank, beta, residual) & $-0.12$ \\
\bottomrule
\end{tabular}
\end{table}

Microstructure features contribute most of the edge.
Removing VPIN alone accounts for $-0.19$ of the microstructure drop.

\subsection{Cost sensitivity}

\begin{table}[h]
\centering
\caption{Sharpe sensitivity to fee and slippage assumptions.}
\label{tab:costs}
\begin{tabular}{lcrr}
\toprule
Fee bps & Slippage mult. & Sharpe & Max DD \\
\midrule
2.0 & 0.5$\times$ & 1.82 & 6.1\,\% \\
3.5 (base) & 1.0$\times$ & 1.41 & 8.2\,\% \\
5.0 & 1.5$\times$ & 1.02 & 10.4\,\% \\
7.0 & 2.0$\times$ & 0.58 & 14.1\,\% \\
\bottomrule
\end{tabular}
\end{table}

The strategy remains marginally profitable (Sharpe $>$ 0.5) up to 2$\times$
the base slippage assumption, providing meaningful cost headroom.

\subsection{Latency sensitivity}

At 15\,ms (co-located) vs.\ 100\,ms (remote VPS) the Sharpe degrades by 0.21.
At 300\,ms the strategy breaks even; at 500\,ms it is unprofitable.
This confirms that co-location or a fast cloud instance (sub-50\,ms to exchange)
is necessary for the directional component.
The funding carry sleeve is latency-insensitive and profitable at any latency.

% =============================================================================
\section{Pitfalls \& Limitations}
\label{sec:pitfalls}

\subsection{Pitfalls found and fixed}

\textbf{Look-ahead via \texttt{shift(0)}.}
An early version of \texttt{OrderFlowImbalance} used the current bar's ask size
directly rather than shifting by one bar.
This leaked future order-book state into the feature vector and inflated
backtest Sharpe by approximately 0.4.
Fix: all features now call \texttt{.shift(1)} before any cross-bar computation,
enforced by the \texttt{point\_in\_time\_safe} flag and Hypothesis property tests.

\textbf{Bar aggregation boundary error.}
When aggregating 1-min bars to 5-min bars, \texttt{resample("5min")} with
\texttt{closed="right", label="right"} was used incorrectly:
the bar labelled at time $T$ included the return from $T$ to $T$+1min,
giving a 1-bar look-ahead.
Fix: switch to \texttt{closed="left", label="right"}; add a regression test
that verifies each bar's close equals the last 1-min close before the label time.

\textbf{Meta-model leakage via primary-model in-sample predictions.}
The meta-model was initially trained on the primary model's in-sample
predictions, not its OOS predictions.
This caused the meta-model to learn the primary model's training-set biases
rather than its generalisation behaviour, inflating meta-model precision
by 8--12\,pp.
Fix: the meta-model is now trained exclusively on OOS predictions from the
same purged folds used to evaluate the primary model.

\textbf{Slippage underestimation during high-volatility bars.}
The square-root impact model was calibrated on average-volatility days.
During the first 30 minutes of a volatility spike (identified by VolOfVol
crossing its 95th percentile), actual slippage is $\approx$\,2$\times$ the
model estimate.
Fix: add a volatility regime multiplier to the slippage model;
reduce position size to 50\,\% when VolOfVol is elevated.

\textbf{HMM state instability at regime boundaries.}
The HMM Viterbi path switches states frequently near regime boundaries,
producing noisy gating signals.
Fix: require the HMM posterior probability to exceed 0.70 for a full bar
before acting on a regime change (one-bar lag confirmation).

\subsection{Limitations}

\begin{enumerate}
  \item \textbf{Data scope}: The backtest covers only Binance BTCUSDT,
        ETHUSDT, and SOLUSDT.
        Results may not generalise to smaller-cap perpetuals with lower
        liquidity and wider spreads.
  \item \textbf{Short live paper run}: 48 hours is insufficient to measure
        live Sharpe with statistical significance.
        A minimum of 30 days is needed for PSR $>$ 0.90 at this return frequency.
  \item \textbf{Order book depth}: All microstructure features that require
        Level-2 data (OFI, MicroPrice, VPIN) fall back to OHLCV approximations
        in backtesting, which understates their live information content.
  \item \textbf{Funding carry sleeve}: The carry sleeve was tested for 3 weeks;
        carry conditions can reverse rapidly in bear markets.
  \item \textbf{Trial count}: The DSR calculation uses $N{=}247$ trials,
        which includes Optuna runs, ablations, and manual notebook experiments.
        Some manual experiments were not logged in the trial registry;
        the true $N$ may be higher, which would lower the DSR further.
\end{enumerate}

% =============================================================================
\section{Future Work}

\begin{enumerate}
  \item \textbf{Tick-level features}: Move from 5-min bar features to
        millisecond-resolution L2 snapshots to improve VPIN and OFI accuracy.
  \item \textbf{Execution RL}: Replace the static limit/market fallback with
        a PPO/SAC agent trained via direct reinforcement~\cite{moody2001rl}
        to optimise slippage on each trade.
  \item \textbf{Longer live run}: Run paper trading for $\ge$\,30 days to
        achieve meaningful PSR estimates and detect live regime drift.
  \item \textbf{Multi-venue arbitrage}: The perpetual basis (SpotPerpBasis)
        sometimes diverges across Binance and Bybit; triangular arbitrage
        across venues would add a low-risk carry sleeve.
  \item \textbf{Adaptive retraining}: Replace fixed fold-boundary retraining
        with online learning triggered by a Bayesian changepoint
        detector~\cite{adams2010changepoint} monitoring feature drift.
  \item \textbf{Broadening the universe}: Test on 20+ symbols to evaluate
        whether cross-sectional features (UniverseRank, BetaToBTC) provide
        additional edge at portfolio level.
\end{enumerate}

% =============================================================================
\section*{Reproducibility}

All experiments are deterministic given the configured random seed
(\texttt{seed\_everything(42)}).
The full pipeline from raw data to figures is reproducible with:

\begin{lstlisting}
make setup
make backtest  # runs full walk-forward, saves results
make figures   # regenerates all paper figures from data
cd paper && pdflatex main.tex && bibtex main && pdflatex main.tex
\end{lstlisting}

Data and model artifacts are excluded from git via \texttt{.gitignore}
but the generation scripts and configs are fully committed.

\bibliographystyle{plainnat}
\bibliography{refs}

\end{document}
